\section{Introduction}
A restoration front end can make a degraded scene easier to see. A detector may still miss objects when restoration or resampling changes their boundaries and texture. On an edge device, restoration also consumes memory and time. Evaluating a cascade therefore requires measurements at the geometry and precision used in each task.

Aegis Perception links these measurements through a shared four-weather benchmark (Fig.~\ref{fig:system}). Six trained systems use a fixed split and report quality, improvement over the degraded input, and parameter count. We include a compact prompt-conditioned LiteWeather v2, four published transformer restorers, and an attention U-Net baseline. A matched continuation control tests whether spatial/detail adapters or auxiliary losses improve on further training of the retained checkpoint.

We evaluate the selected Restormer with a frozen YOLOv8n at distinct operating points. An earlier local pilot includes a resize-only control. Labelled Nano replay measures the deployed $128\times128$ cascade, and live and saved-camera workloads characterize its execution cost. A separate local diagnostic holds source pixels and the detector fixed across six paths to examine the effect of resolution. A per-channel weight-grid screen probes the sensitivity of convolutional attention projections and feed-forward layers before attempting low-bit execution.
